\clearpage\thispagestyle{otFront}\phantomsection\label{parcours}\pdfbookmark[0]{Construire son parcours}{bm-parcours}\otPageTitre{Construire son parcours}{Une année avec des retours réguliers}\fontsize{11}{15}\selectfont

\begin{tabularx}{\linewidth}{@{}lY@{}}\toprule
Calculer & 1 à 5 : grands nombres, opérations, division, égalités et suites.\\[3mm]
Partager & 6 à 10 : fractions, fractions décimales et nombres décimaux.\\[3mm]
Mesurer & 11 à 14 : longueurs, masses, contenances, aires, angles et durées.\\[3mm]
Construire & 15 à 18 : droites, cercles, figures usuelles, symétrie et solides.\\[3mm]
Modéliser & 19 à 21 : proportionnalité, données et premiers pas vers le hasard.\\\bottomrule
\end{tabularx}
\begin{savoirfaire}[Une séance qui fonctionne]Relire un exemple du cours,
chercher deux exercices sans corrigé, comparer, puis refaire la question
manquée. Dix à quinze minutes régulières valent mieux qu’une longue
séance isolée.\end{savoirfaire}
\begin{savoirfaire}[Revenir et mélanger]Reprendre une notion quelques jours
plus tard, puis au moment des bilans. Les problèmes mélangent
volontairement les chapitres : c’est ainsi qu’on apprend à choisir.\end{savoirfaire}
Les défis demandent entre 20 et 45 minutes, parfois en plusieurs fois.
Le tableau de bord suit les réussites et les reprises au fil de l’année.

\clearpage
